\section{Credit assignment across tools and memory}
\label{sec:C08:credit}
A terminal reward must assign credit to decisions separated by long tool delays and many intermediate tokens. Treating every token in a successful trajectory as equally good can reinforce redundant explanations and accidental workarounds. Tool selection, argument construction, context editing, and memory writing have different causal roles. A useful ablation removes or replaces one component while preserving resources and observes the outcome. Such interventions can inform training signals, but hindsight attribution remains uncertain when several actions jointly enable success.

Memory writing is particularly difficult because its value may appear only in a later session. Immediate task reward can favor a compact note that helps the current answer while discarding a rare future constraint. Conversely, rewarding every retrieved use can encourage excessive storage or irrelevant retrieval. The learning objective should declare the horizon and cost of retention, retrieval, verification, and conflict handling. Evaluation should randomize task order and distinguish within-session adaptation from cross-session transfer. Otherwise improvements can depend on a favorable curriculum that deployment will not reproduce.

Model updates also change the policy that interprets existing memories and skills. A procedure distilled for one checkpoint may become redundant, misleading, or incompatible after post-training. Evaluate retained skill banks before and after a weight update, including stale assumptions and changed tool schemas. Updating memory and model weights together without separate ablations makes it difficult to identify where a gain or regression originated. Compatibility tests should therefore be part of the learning loop, with rollback capable of restoring the relevant combination of policy and external state.

Post-training can improve reasoning and tool use substantially, but the claim must remain attached to a measurable behavior and evidence source. Preference agreement is not universal alignment, a higher automated reward is not verified task success, and a finite sequence of improving checkpoints is not accelerating recursive self-improvement. The scientific standard is stronger: correct conditioning, identified estimators, independent acceptance criteria, resource-matched alternatives, and retained regressions. These requirements make the resulting policy improvement interpretable and useful for a stateful agent rather than merely attractive on its training dashboard.
